% Fragmento conservado; véase MANIFIESTO_ANTECEDENTES_LECTURA.json.
\section{Dodecaphase record and the rational phase constant}
\label{sec:k-registro}

Generation of the regional coordinates does not exhaust the information in the
APP--TRIT--TPK state. One Archimedean readout determines a value through
compatible cylinders; another readout preserves the opposition of the sheets,
the orientation of the route, and the record of the incidences traversed. The
dodecaphase record belongs to this second readout before entering
the closure of alpha. It is therefore not defined by subtracting a physical constant from
a combination of previously known values. It is reconstructed from a
signed coordinate of the same arithmetic state with trajectory memory that produces the regional
coordinates.

The distinction is necessary even when two objects have twelve
coordinates. A twelve-block catalogue word, a list of twelve
oriented events, a signed integer vector, and a base-thousand dodecaphase record are
different objects. Comparing them requires the maps that
relate them. Equality of length does not provide these maps, and
loss of signs is not merely a simplification of notation.

This chapter preserves the causal cutoff established in the preceding chapters:
the nine APP positions on the horizontal and vertical axes produce the
arithmetic sheets; TRIT types regime and orientation; TPK selects,
transports, updates, and prolongs, preserving residue, quotient, carry,
sheet, route, boundary, and memory. The coordinates used here are obtained
from the joint discrete structure of the continuum. The chain
\(w_6\to w_{12}\to w_{18}\to w_{24}\to w_{30}\to R_{36}\to G_9\)
is not replaced by a finite schedule. Nor is the order
\(U_t\to\Gamma_9\to K_{\mathrm{ph}}\) reversed: the last map is a
reduced-memory readout, not the generator of the signed record.

% Preserved fragment; see MANIFIESTO_ANTECEDENTES_LECTURA.json.
% Adición propuesta al inicio de registro_k.tex, después de su introducción.
% Conserva todas las definiciones, fórmulas y demostraciones actuales.
% Procedencia: registro_k.tex, §§k-hadamard/k-transversal/k-kappas;
% k_reversibilidad.tex; incidencia_articulacion.tex, inc-ampl-cuaterna.
\paragraph{Reversible coordinatization and incidence determination.}
\label{ampl-alfa-k-intro-reversible}
The result organizing this chapter is the equivalence of three
coordinatizations of a register previously produced by
APP--TRIT--TPK. The Hadamard transformation relates the signed register
to the integer register; the difference extractor relates
the latter to its transverse variations and its uniform component:
\[
 U\underset{\mathcal H_{12}}{\overset{\mathcal H_{12}/4}{\longleftrightarrow}}K
 \overset{(D_3,D_4,Q)}{\longleftrightarrow}
 (D_3K,D_4K,Q(K)).
\]
The first equivalence is established on the integral sublattice
$\mathcal L_H$; the second, on the image of the extractor. With
base $1000$, the $12$ positions and the phase origin fixed, the evaluation
$K\mapsto\kappa_{\mathrm{per}}$ adds a fourth reversible
coordinatization: multiplying by $1000^{12}-1$ and performing the
Euclidean divisions recovers the full register. The proofs in
\S\ref{subsec:k-hadamard}, \S\ref{subsec:k-transversal} and
\S\ref{subsec:k-kappas} establish these equivalences in their respective
domains.

Two compositions subsequently act on this register. The positional
composition combines its components with the closure,
propagation and self-scaling coordinates, and its normalization determines $\alpha_A$. The
incidence composition applies the cubic threshold and yields
\[
 K\longmapsto B_K=\{j:K_j\geq729\}=\{4,6,9,10\}.
\]
The full register allows both compositions to be reconstructed. The support
$B_K$ specifically retains membership in the class selected
by the threshold. The identity $B_K=H_e\cap P_\pi$ in
(8.12) of Article I (provenance of the exceptional prolongation, preserved in the technical record) connects this integer extraction to the incidence
of the regional codes. The negative support of the balance with $K$
completes the flag $B_K\subset H^-_\alpha$ developed in the
exceptional prolongation. Thus the role of $K$ is defined by
an invertible transformation and by precise subsequent
maps; the additional information on depth and transport remains
in the state from which the register originates.


\subsection{Cyclic extension and memory grading}
\label{subsec:k-calendario}

The observable schedule contains one hundred and eight positions, organized into
twelve nonadic windows. With positive indices, let
\[
 E_m=\{9(m-1)+1,\ldots,9m\},\qquad 1\leq m\leq12.
\]
The ordered support \(E_1\sqcup\cdots\sqcup E_{12}\) preserves position and
window membership. The group \(C_{108}=\mathbb Z/108\mathbb Z\),
by contrast, preserves the law of cyclic translation. For the representative
\(0\leq\widetilde\theta<108\), the coordinates
\[
 \beta(\theta)=\left[\left\lfloor\frac{\widetilde\theta}{9}\right\rfloor\right]_{12},
 \qquad r(\theta)=1+(\widetilde\theta\bmod9)
\]
separate the dodecaphase sector from the interior nonadic position. Neither
by itself preserves the total number of turns.

\begin{proposition}[The carry of the schedule]
\label{prop:k-extension}
The maps
\[
 \iota:C_{12}\longrightarrow C_{108},\quad[b]\longmapsto[9b],
 \qquad p:C_{108}\longrightarrow C_9,\quad[\theta]\longmapsto[\theta]_9
\]
form a nonsplit exact sequence
\[
 0\longrightarrow C_{12}\xrightarrow{\iota}C_{108}
 \xrightarrow{p}C_9\longrightarrow0.
\]
For the set-theoretic section choosing the representatives
\(0,\ldots,8\), the failure of additivity is the carry
\[
 c(r,r')=\left[\left\lfloor\frac{\widetilde r+\widetilde r'}9\right\rfloor\right]_{12}.
\]
\end{proposition}

\begin{proof}
The kernel of \(p\) consists of the multiples of nine and coincides with
the image of \(\iota\). Reduction is surjective, and multiplication
by nine is injective on the stated domain. If the sequence
split, the middle group would be isomorphic to \(C_{12}\times C_9\), whose
exponent is \(\operatorname{mcm}(12,9)=36\); the exponent of
\(C_{108}\) is 108. This contradiction proves nonsplitting. Finally,
Euclidean division of \(\widetilde r+\widetilde r'\) by nine gives the
stated carry and verifies the identity for the section.
\end{proof}

The result explains a distinction that will reappear in alpha: returning to the
phase is not equivalent to returning to the state. If \(q_{\mathrm{mem}}\) counts
nonadic turns, the schedule preserves only
\(\beta=[q_{\mathrm{mem}}]_{12}\). After one hundred and eight steps,
\(\beta\) returns, but \(q_{\mathrm{mem}}\) increases by twelve. The
record also preserves the cylinder, the incidence ledger, and the vertical
memory. Closure of a finite chart does not turn this history into a
memoryless cycle.

\subsection{The oriented record and the integral \texorpdfstring{\(1+5+6\)}{1+5+6} chart}
\label{subsec:k-registro-integral}

For a history produced by TPK, let \(\tau_m\) be the subroute over
\(E_m\), \(\sigma_m\) its orientation, \(\kappa_m\) the boundary
carry, and \(\Lambda_m\) the local incidence ledger. The record is
\[
 \mathcal R_{12}(x)=
 \bigl((E_m,\tau_m,\sigma_m,\kappa_m,\Lambda_m)\bigr)_{m=1}^{12}.
\]
Its domain is the space of arithmetic histories with trajectory memory, not the set of
phase labels. Projection onto the schedule forgets precisely some
of the data required by the signed readout.

An initial chart of the record can be written in terms of oriented
events. The reference construction~\textup{(provenance preserved in the technical register)} fixes the event code set
\(\{2,4,6,8\}\) and the sector orientation
\[
 \Sigma_{12}=(+,+,+,-,-,-,+,+,+,-,-,-).
\]
If \(d_t\) is the direction code at step \(t\), this gives
\[
 e_i=\Sigma_{12}(i+1)
 \#\{t\in E_{i+1}:d_t\in\{2,4,6,8\}\},
 \qquad0\leq i<12.
\]
The formula keeps the nonnegative event count separate from
orientation. It does not identify a cardinal direction with a TRIT sign. The
vector \(e\), whose components are signed local counts, is also not
yet the vector \(U\) that will appear below.

Opposite positions determine, for \(0\leq i<6\),
\[
 p_i=e_i+e_{i+6},\qquad a_i=e_i-e_{i+6}.
\]
With
\[
 s_C=\sum_{i=0}^{5}p_i,\qquad M_i=p_i-p_5\quad(0\leq i<5),
\]
the integral chart is defined by
\[
 \mathcal L_{12}(e)
 =(s_C,M_0,\ldots,M_4,a_0,\ldots,a_5).
\]
The decomposition \(1+5+6\) is not an arbitrary allocation of coordinates:
one coordinate preserves the charge, five compare the pairs with a
reference pair, and six preserve the opposition of the half-crowns.

\begin{lemma}[Inversion of the integral chart]
\label{lem:k-carta-integral}
An integer tuple \((s_C,M_0,\ldots,M_4,a_0,\ldots,a_5)\) belongs to the
image of \(\mathcal L_{12}\) if and only if
\[
 s_C-\sum_{i=0}^{4}M_i\equiv0\pmod6,
 \qquad p_i\equiv a_i\pmod2\quad(0\leq i<6),
\]
where
\[
 p_5=\frac{s_C-\sum_{i=0}^{4}M_i}{6},\qquad p_i=p_5+M_i\quad(i<5).
\]
In that case its preimage is unique and is given by
\[
 e_i=\frac{p_i+a_i}{2},\qquad e_{i+6}=\frac{p_i-a_i}{2}.
\]
\end{lemma}

\begin{proof}
The sum of the five margins is \(s_C-6p_5\), which gives the
division by six. Once the \(p_i\) are reconstructed, each pair of equations
\(p_i=e_i+e_{i+6}\), \(a_i=e_i-e_{i+6}\) has the indicated solution.
This solution is integral exactly when the parity congruence holds. All
operations are inverted, so no residual freedom remains.
\end{proof}

This reversibility shows what information is preserved; it does not authorize
discarding the rest of the ledger. The complete record has a larger domain
than this count chart. Extraction of the channels feeding
\(U\) uses that complete record and cannot be deduced from the list
\((e_i)\) by identifying names.

\subsection{Two readouts of the terminal state}
\label{subsec:k-terminal}

Let \(x_{\mathrm{term}}\) denote the state produced by the composition
of selection, transport, and update up to the terminal cutoff. The
operator expression is
\[
 x_{\mathrm{term}}=
 (\mathcal U_{108}\circ\cdots\circ\mathcal U_1)(x_{\mathrm{seed}}),
 \qquad \mathcal U_t=\mathsf{Upd}_t\circ\mathsf{Tra}_t\circ\mathsf{Sel}_t.
\]
Its two relevant readouts are
\[
 x_{\mathrm{term}}
 \xrightarrow{(\pi_{\mathrm{term}},R_{\mathrm{sgn}})}
 (B_{\mathrm{term}},U).
\]
The matrix \(B_{\mathrm{term}}\) retains the visible ternary readout; the
vector \(U\) preserves another coordinate of the same history. No arrow
\(B_{\mathrm{term}}\to U\) is inserted between them.

The terminal selector of the reference construction~\textup{(provenance preserved in the technical register)} acts on catalogues
of sizes 196, 197, and 196. The local signatures reduce their
\(196\cdot197\cdot196=7\,567\,952\) triples to 331, and the column
margin \(q_{\partial}=(1,2,2,1,2,0)\) leaves two matrices:
\[
 B_0=\begin{pmatrix}0&2&0&2&2&0\\0&0&1&2&2&0\\1&0&1&0&1&0\end{pmatrix},
 \qquad
 B_1=\begin{pmatrix}0&2&1&0&2&0\\0&0&1&2&2&0\\1&0&0&2&1&0\end{pmatrix}.
\]
The row orientation is \(\sigma=(1,1,-1)=(1,1,2)\) in
\(\mathbb F_3^3\). The row charges of the candidates are
\((0,2,0)\) and \((2,2,1)\). Only the second belongs to
\(\langle\sigma\rangle=\{(0,0,0),(1,1,2),(2,2,1)\}\); hence
\(B_{\mathrm{term}}=B_1\). This check proves the final step of the
selector from the stated census. It does not turn the column margin
into a consequence of the row charge, nor replace the preceding enumeration
with inspection of these two matrices.

The signed readout has the typed composition
\begin{equation}
 R_{\mathrm{sgn}}
 =\Pi_H\circ\mathcal E_{108}^{90,120}\circ\mathcal R_{12}:
 \mathcal X_{\mathrm{TPK}}^{\mathrm{term,mem}}
 \longrightarrow\mathbb Z^{12}.
 \label{eq:k-lector-firmado}
\end{equation}
The channel record determines
\[
 \mathcal E_{108}^{90,120}(\mathcal R_{12}(x_{\mathrm{term}}))
 =(b^{90},b^{120},Q_{\mathrm{TPK}}),
\]
with coordinates
\begin{align*}
 b^{90}={}&(-495,-116,-684,108,601,-90,\\
            &\hspace{19mm}-173,-88,313,560,-397,461),\\
 b^{120}={}&(-425,-281,-481,671,-255,30,\\
             &\hspace{19mm}475,-543,680,251,6,-128),\\
 Q_{\mathrm{TPK}}={}&6263.
\end{align*}
These quantities are used as coordinates of the preceding record, not
as parameters fitted through alpha. The reconstruction proved
in the following subsections begins exactly at this typed output.
The material trace of the preceding extractor and the scope of its
self-contained transcription are preserved in the technical provenance note; the subsequent
arithmetic is not presented as a substitute for that trace.

\subsection{From the channel record to the signed vector}
\label{subsec:k-pi-h}

The harmonic readout \(\Pi_H\) can be written without using either \(K\) or
alpha. To avoid confusing its orbit sums with the alpha blocks,
we denote them by \(s_1,s_2,s_3\). Let
\[
 B_r=\sum_{j=0}^{3}b^{120}_{r+3j},\qquad
 d_{r,j}=b^{90}_{r+3j},\qquad1\leq r\leq3,
\]
and
\begin{equation}
 s_1=\frac{Q_{\mathrm{TPK}}+2B_1+B_2}{3},\qquad
 s_2=s_1-B_1,\qquad s_3=s_2-B_2.
 \label{eq:k-sumas-armonicas}
\end{equation}
Division by three belongs to the readout of the three orbits. Its
integrality is checked on compatible outputs of the record; it is not
assumed for an arbitrary element of \(\mathbb Z^{25}\).

The three signed blocks are
\begin{equation}
 U^{(r)}=(s_r,d_{r,1}-d_{r,3},d_{r,0}+d_{r,2},d_{r,0}-d_{r,2}).
 \label{eq:k-bloques-firmados}
\end{equation}
Integer substitution gives
\[
 (B_1,B_2,B_3)=(972,-1073,101),\qquad
 (s_1,s_2,s_3)=(2378,1406,2479),
\]
and therefore
\begin{align*}
 U^{(1)}&=(2378,-452,-668,-322),\\
 U^{(2)}&=(1406,998,-204,-28),\\
 U^{(3)}&=(2479,-551,-371,-997).
\end{align*}
Interleaving by columns yields
\begin{equation}
 U=(2378,1406,2479,-452,998,-551,
       -668,-204,-371,-322,-28,-997).
 \label{eq:k-u-firmado}
\end{equation}
The presence of negative values and components greater than 999
shows that \(U\) is not a base-thousand digit word. Nor is it
\(U_6\Vert U_6\), a catalogue concatenation of period six. The distinction
concerns domain and preserved information, not merely values.

The constructive order thus far is
\[
 x_{\mathrm{term}}\longrightarrow\mathcal R_{12}
 \longrightarrow(b^{90},b^{120},Q_{\mathrm{TPK}})
 \longrightarrow U.
\]
The inversion that follows produces the dodecaphase record. The fact that the
channels can subsequently be recovered from the dodecaphase record is a reversibility check; it does not
reverse this causal order.

\subsection{Orbitwise Hadamard transformation and integer reconstruction}
\label{subsec:k-hadamard}

In the twelve-position cycle, step three determines three orbits:
\[
 (1,4,7,10),\qquad(2,5,8,11),\qquad(3,6,9,12).
\]
The separation of four sectors in each orbit is what, once
origin and orientation have been fixed, admits the ninety-degree angular readout. The
transformation does not act on three consecutive tetrads in the natural ordering.

Let
\[
 H_4=\begin{pmatrix}
 1&1&1&1\\1&1&-1&-1\\1&-1&1&-1\\1&-1&-1&1
 \end{pmatrix}.
\]
If \(P_{\mathrm{orb}}\) reorders by the three indicated orbits,
define
\[
 \mathcal H_{12}=P_{\mathrm{orb}}^{-1}
 (H_4\oplus H_4\oplus H_4)P_{\mathrm{orb}}.
\]

\begin{lemma}[Hadamard inversion and integrality]
\label{lem:k-hadamard}
We have
\[
 H_4^2=4I_4,\qquad H_4^{-1}=\frac14H_4,\qquad\det H_4=-16.
\]
For \(u\in\mathbb Z^4\), the preimage \(H_4^{-1}u\) is integral if and
only if \(H_4u\equiv0\pmod4\).
\end{lemma}

\begin{proof}
The rows are orthogonal and have squared norm four; the matrix is
symmetric. This gives \(H_4^2=4I_4\) and the inverse. Integer elimination,
or expansion of the determinant, gives the negative sign and the value \(-16\).
The integrality condition is exactly divisibility by four
of all four coordinates of \(H_4u\).
\end{proof}

\begin{definition}[Integral dodecaphase transformation]
The sublattice of signed records and its integral coordinatization are
\[
 \mathcal L_H=\{u\in\ZZ^{12}:\mathcal H_{12}u\equiv0\pmod4\},
 \qquad \mathscr K(u)=\tfrac14\mathcal H_{12}u.
\]
The map \(\mathscr K:\mathcal L_H\to\ZZ^{12}\) is an isomorphism
of abelian groups, with inverse \(k\mapsto\mathcal H_{12}k\).
The canonical record \(K\) is the image of the signed record generated
by the preceding operations, with phase origin and orientation fixed.
\end{definition}

\begin{proof}
The congruence defines precisely the integrality of \(\mathscr K(u)\).
The identity \(\mathcal H_{12}^2=4I\) proves that both compositions
are inverse. In particular, every \(k\in\ZZ^{12}\) has preimage
\(\mathcal H_{12}k\in\mathcal L_H\).
\end{proof}

\begin{proposition}[Canonical integral record]
\label{prop:k-sello}
The inversion \(K^{(r)}=H_4U^{(r)}/4\) of the blocks in
\eqref{eq:k-bloques-firmados} produces
\begin{align*}
 K^{(1)}&=(234,729,621,794),\\
 K^{(2)}&=(543,659,58,146),\\
 K^{(3)}&=(140,824,914,601).
\end{align*}
In the natural ordering, the dodecaphase record is
\begin{equation}
 K=(234,543,140,729,659,824,621,058,914,794,146,601).
 \label{eq:k-vector}
\end{equation}
It is the unique rational preimage of \(U\) and belongs to
\(\{0,\ldots,999\}^{12}\).
\end{proposition}

\begin{proof}
The three products before division are
\begin{align*}
 H_4U^{(1)}&=(936,2916,2484,3176),\\
 H_4U^{(2)}&=(2172,2636,232,584),\\
 H_4U^{(3)}&=(560,3296,3656,2404).
\end{align*}
All coordinates are divisible by four. Dividing and undoing the
orbit permutation gives \eqref{eq:k-vector}. Uniqueness follows from
rational invertibility, not from a search among candidates close to
a physical value.
\end{proof}

Integrality can also be formulated in lattice terms. The
determinantal divisors of \(H_4\) are \(1,2,4,16\): the entries
have greatest common divisor one; the minors of order two are even and
one has absolute value two; those of order three are multiples of four and
one has absolute value four. Hence
\[
 \operatorname{SNF}(H_4)=\operatorname{diag}(1,2,2,4),
\]
and for the global transformation,
\[
 \operatorname{coker}\mathcal H_{12}
 \cong(\mathbb Z/2\mathbb Z)^6\oplus(\mathbb Z/4\mathbb Z)^3.
\]
The index of the image is \(16^3=4096\). Thus, not every signed
integer vector admits an integer preimage; the vector obtained satisfies the
precise congruences required by inversion. The factor \(1/4\) is
forced by the matrix and does not have the status of a free coefficient.

\subsection{Cyclic differences and the uniform component}
\label{subsec:k-transversal}

With cyclic indices modulo twelve, let
\[
 (D_3z)_m=z_m-z_{m+3},\qquad(D_4z)_m=z_m-z_{m+4},
 \qquad Q(z)=\sum_{m=1}^{12}z_m.
\]
The first operator compares sectors separated by three positions; the
second, by four. The channel names 90 and 120 refer here to
these separations in the oriented cycle, not to identification with a
physical operator from another domain.

\begin{proposition}[Completeness of the transverse system]
\label{prop:k-transversal}
The operator
\[
 \mathcal D=(D_3,D_4,Q):\mathbb Z^{12}\longrightarrow\mathbb Z^{25}
\]
has rank twelve and nonzero invariant factors
\[
 \operatorname{diag}(1,\ldots,1,12),
\]
with eleven units. In particular, its outputs determine a unique dodecaphase record.
\end{proposition}

\begin{proof}
If \(D_3z=D_4z=0\), the vector is constant along both steps.
Since \(3\) and \(4\) generate the cycle modulo twelve,
\(\ker D_3\cap\ker D_4=\langle\mathbf1\rangle\). The total charge
eliminates this last freedom because \(Q(\mathbf1)=12\).

For the integral statement, the rows of \((D_3,D_4)\) are oriented incidences
of a connected graph. A spanning tree provides eleven
unit invariant factors. Every minor of order twelve must contain the
row \(Q\). Adding the first eleven columns to the last makes the latter
zero in the incidence rows and twelve in row \(Q\). All
maximal minors are divisible by twelve, and the tree produces one
of absolute value exactly twelve. The final factor is therefore twelve.
\end{proof}

Checking \eqref{eq:k-vector} gives
\[
 D_3K=b^{90},\qquad D_4K=b^{120},\qquad Q(K)=6263.
\]
Formulas \eqref{eq:k-sumas-armonicas} and
\eqref{eq:k-bloques-firmados}, in turn, reconstruct \(U\) from those
differences. This closes a triangle of exact representations of the
same object: transverse record, signed vector, and dodecaphase record. Equality
of results does not turn their codomains into a single one; it proves that
the specified maps recover the relevant information.

\subsection{Two rational readouts: \texorpdfstring{\(\kappa_{36}\)}{kappa36} and \texorpdfstring{\(\kappa_{\mathrm{per}}\)}{periodic kappa}}
\label{subsec:k-kappas}

Once \(K\) has been constructed, the base-thousand chart allows two different
operations. The first reads a single turn; the second prolongs the dodecaphase record
periodically. Let \(B=1000\) and define its concatenation integer
\begin{equation}
 N_K=\sum_{m=1}^{12}K_mB^{12-m}
 =234543140729659824621058914794146601.
 \label{eq:k-numerador}
\end{equation}
The leading zeros of a block, as in 058, are part of its three-digit
representation. They do not change its integer value, but they make it possible to concatenate
the twelve blocks unambiguously.

\begin{definition}[Finite and periodic readouts of the dodecaphase record]
\label{def:k-kappas}
The one-turn readout is
\[
 \kappa_{36}=\sum_{m=1}^{12}K_mB^{-m}=\frac{N_K}{B^{12}}.
\]
The periodic prolongation is
\[
 K_j^{\mathrm{per}}=K_{1+((j-1)\bmod12)},\qquad
 \kappa_{\mathrm{per}}=\sum_{j\geq1}K_j^{\mathrm{per}}B^{-j}.
\]
\end{definition}

\begin{proposition}[Exact value and period of the prolonged dodecaphase record]
\label{prop:k-periodo}
We have
\begin{equation}
 \kappa_{\mathrm{per}}=\frac{N_K}{B^{12}-1},\qquad
 \kappa_{\mathrm{per}}-\kappa_{36}
 =\frac{N_K}{B^{12}(B^{12}-1)}
 =B^{-12}\kappa_{\mathrm{per}}.
 \label{eq:k-diferencia-kappas}
\end{equation}
The infinite word of the dodecaphase record has minimum period twelve in base thousand and
minimum period thirty-six in base ten.
\end{proposition}

\begin{proof}
The sum of the first turn is \(N_KB^{-12}\); each new turn
multiplies its contribution by \(B^{-12}\). The geometric series gives
\(N_K/(B^{12}-1)\), and subtraction gives the second identity.

The twelve components of \(K\) are distinct. A proper period of the cycle
would identify two of them; the minimum block period is therefore
twelve. The decimal expansion is canonical, since the dodecaphase record has no tail
of blocks equal to 999. If its minimum decimal period were \(d\), it would divide
36. Grouping digits in threes would produce a block period of
\(d/\gcd(d,3)\). No proper divisor of 36 gives twelve under this
operation. Hence the minimum decimal period is 36.
\end{proof}

Both readouts are rational, but they are not the same rational number. The first
terminates after twelve blocks; the second repeats the turn. Moreover,
\(0<\kappa_{\mathrm{per}}-\kappa_{36}<B^{-12}\). Their proximity does not
allow one to be substituted for the other in an exact identity. In particular,
the finite closure of alpha uses \(\kappa_{36}\), whereas the
sum of the prolonged section uses \(\kappa_{\mathrm{per}}\).

% Fragmento conservado; véase MANIFIESTO_ANTECEDENTES_LECTURA.json.
\subsection{Reversibility of the rational constant and phase change}

The rational constant \(\kappa_{\mathrm{per}}\) completely determines
the register \(K\) when the base, length and readout origin
are preserved. Its numerical and structural functions are related
through an exact inversion.

\begin{proposition}[Integral recovery and cyclic transformation]
Let \(I(K)=\sum_{j=1}^{12}K_j1000^{12-j}\). Then
\[
 I(K)=(1000^{12}-1)\kappa_{\mathrm{per}}.
\]
The integer expansion of \(I(K)\) in twelve base-\(1000\) blocks
recovers all \(K_j\), including the leading zeros of each block.
If \(\rho(K)=(K_2,\ldots,K_{12},K_1)\), one has
\begin{equation}
 \kappa(\rho K)=1000\kappa(K)-K_1.
 \label{eq:k-rotacion-racional}
\end{equation}
\end{proposition}
\begin{proof}
The first identity inverts the definition of periodic evaluation.
Iterated Euclidean division by \(1000\), with twelve fixed positions,
recovers the vector. Moreover,
\[
 I(\rho K)=1000I(K)-K_1(1000^{12}-1).
\]
Dividing by \(1000^{12}-1\) gives \eqref{eq:k-rotacion-racional}.
\end{proof}

Recovery therefore composes \(\kappa\mapsto K\mapsto U\)
and the transformations \(D_3K,D_4K,Q\) already constructed. Rational
evaluation preserves the signed register and the cyclic differences
depending on it. A history's total depth and transport operators
remain in their own coordinates: the preceding inversion
has the integral dodecaphasic register as its domain.

Write \(\kappa_j=\kappa(\rho^jK)\), with indices modulo \(12\).
The same identity gives
\[
 K_{j+1}=1000\kappa_j-\kappa_{j+1},\qquad
 \sum_{j=0}^{11}\kappa_j=\frac{\sum_{j=1}^{12}K_j}{999}
 =\frac{6263}{999}.
\]
Phase acts exactly on the constant: its value is fixed
in the canonical readout, while its rotations follow an affine law.
The orbit sum is invariant under choice of origin.

Subsequent incidence also has an expression in these
coordinates. The cubic threshold determines
\[
 B_K=\{j+1:1000\kappa_j-\kappa_{j+1}\ge729\}
 =\{4,6,9,10\}.
\]
The tetrad is recovered from the rational orbit because that orbit
first recovers the integer components. The signed support associated
with \(\alpha\) additionally uses its regional registers and boundary
normalization. The relation between coordinate and incidence is thus
expressed through identifiable operations.

This is an operational definition of the function of \(K\): it provides
a reversible integral coordinatization of the oriented register, whose
periodic evaluation and incidence functions participate in
subsequent constructions. Value, transformation and domain
are preserved jointly.


\subsection{Periodicity and the cokernel of the transport operator}
\label{subsec:k-clase-acarreo}

The periodicity just proved belongs to the dodecaphase record. It does not imply
periodicity of the regional coordinates, of the closure carries,
or of alpha. It is useful to make this separation precise through a second,
purely integral quotient.

Let \(S\) be the cyclic shift of twelve coordinates, with fixed
orientation, and let \(b\geq2\) be an integer base. The periodic carry
operator is \(\partial_b=S-bI\). If a periodic identity is written
\[
 A=P+E-\Phi-K+\partial_bC,
\]
then the transformations
\[
 K\longmapsto K+\partial_bh,\qquad C\longmapsto C+h
\]
preserve its right-hand side. This invariance compares representatives of a
carry class; it is not authorization to modify the dodecaphase record already
selected by the record.

\begin{proposition}[Periodic carry quotient]
\label{prop:k-cociente-acarreo}
With the preceding orientation,
\[
 \operatorname{coker}(S-bI)\cong\mathbb Z/(b^{12}-1)\mathbb Z.
\]
\end{proposition}

\begin{proof}
The cyclic module is presented as
\(\mathbb Z[t]/(t^{12}-1)\). Imposing \(S=bI\) amounts to imposing
\(t=b\), leaving the relation \(b^{12}-1=0\). Evaluating the
coefficients at \(b\), with the chosen concatenation order, represents
the resulting class.
\end{proof}

The denominator \(B^{12}-1\) of \(\kappa_{\mathrm{per}}\) and this
quotient come from the same cyclic length, but their objects remain
different: one is a rational evaluation and the other a quotient
of a carry module. The finite closure of alpha will be an oriented
chain with a right boundary. Its prolongation will transport a
compatible remote boundary, not a carry identified by decree
between the first and thirteenth positions.
